\chapter{Analyse et conception}
\label{ch:analyse-conception}

\section{Acteurs et cas d'utilisation}

Le système distingue cinq rôles internes, organisés en une hiérarchie de
sécurité (\texttt{ROLE\_SUPER\_ADMIN > ROLE\_ADMIN}), à laquelle s'ajoutent
deux acteurs externes ne possédant pas de compte : le tuteur de stage
(accès par lien à usage unique) et le visiteur public (formulaire de
contact, infolettre, vérification de certificat).

\begin{table}[H]
  \centering
  \caption{Rôles applicatifs et responsabilités}
  \label{tab:roles}
  \begin{tabular}{@{}p{2.8cm}p{9.5cm}@{}}
    \toprule
    Rôle & Responsabilités \\
    \midrule
    \texttt{SUPER\_ADMIN} & Paramètres système globaux, en plus de tous les droits Admin. \\
    \texttt{ADMIN} (Directeur) & Gestion pédagogique complète, validation des UF sensibles, gestion des groupes et des campagnes email, en plus de tous les droits Formateur. \\
    \texttt{FORMATEUR} & Studio de cours, gestion des quiz, correction manuelle, messagerie, bourse à l'emploi. \\
    \texttt{ETUDIANT} & Espace apprenant : contenu, quiz, stage, certificat, messagerie. \\
    \texttt{SUPPORT} & Accès limité, essentiellement en lecture. \\
    \bottomrule
  \end{tabular}
\end{table}

La Figure~\ref{fig:cas-utilisation} synthétise l'ensemble des cas
d'utilisation identifiés à partir des 29 contrôleurs REST exposés par le
backend.

\begin{figure}[H]
  \centering
  \includegraphics[width=\textwidth,height=0.85\textheight,keepaspectratio]{diagramme_cas_utilisation.png}
  \caption{Diagramme de cas d'utilisation}
  \label{fig:cas-utilisation}
\end{figure}

\section{Modèle de données}
\label{sec:analyse-conception}

Le modèle de données regroupe 39 entités JPA concrètes, toutes héritant
--- à l'exception d'\texttt{AssetDownload} --- d'une superclasse technique
abstraite \texttt{AuditableEntity} (\texttt{createdAt}/\texttt{updatedAt}).
La Figure~\ref{fig:classes-ensemble}
en donne une vue d'ensemble au format UML standard, organisée en huit
paquetages fonctionnels avec leurs attributs principaux, leurs
associations et les cardinalités correspondantes. Trois domaines jugés
les plus structurants pour la compréhension du système sont ensuite
détaillés à un niveau plus fin : structure pédagogique et groupes
(Figure~\ref{fig:classes-pedagogie}), évaluation
(Figure~\ref{fig:classes-evaluation}), et stage et certification
(Figure~\ref{fig:classes-stage}). Les domaines restants (communication,
marketing/emploi, médias, suivi système) suivent les mêmes conventions de
modélisation et sont détaillés au besoin dans le code source
(\texttt{domain/entity/}).

\begin{sidewaysfigure}
  \centering
  \includegraphics[width=\textwidth,height=0.9\textheight,keepaspectratio]{diagramme_classes_ensemble.png}
  \caption{Diagramme de classes — vue d'ensemble de l'application}
  \label{fig:classes-ensemble}
\end{sidewaysfigure}

\begin{figure}[H]
  \centering
  \includegraphics[width=\textwidth,height=0.85\textheight,keepaspectratio]{diagramme_classes_pedagogie.png}
  \caption{Diagramme de classes — structure pédagogique et groupes}
  \label{fig:classes-pedagogie}
\end{figure}

\begin{figure}[H]
  \centering
  \includegraphics[width=\textwidth,height=0.85\textheight,keepaspectratio]{diagramme_classes_evaluation.png}
  \caption{Diagramme de classes — évaluation}
  \label{fig:classes-evaluation}
\end{figure}

\begin{figure}[H]
  \centering
  \includegraphics[width=\textwidth,height=0.85\textheight,keepaspectratio]{diagramme_classes_stage_certification.png}
  \caption{Diagramme de classes — stage et certification}
  \label{fig:classes-stage}
\end{figure}

\subsection{Un choix de modélisation délibéré : l'UF comme attribut}

Un choix de conception mérite d'être souligné : l'Unité de Formation (UF)
et l'année du cursus ne sont \emph{pas} modélisées comme des entités
séparées avec leur propre table et clé primaire. Ce sont des attributs
(\texttt{ufCode}, \texttt{ufTitle}, \texttt{yearNumber}) portés directement
par \texttt{ModuleEntity}. Ce choix évite une jointure supplémentaire pour
une notion qui, dans les règles métier réelles, sert uniquement de
\emph{clé de regroupement} (voir \S\ref{sec:progression-uf}) --- il garde
le schéma relationnel simple sans sacrifier la logique de progression par
UF, qui est recalculée dynamiquement à partir de ces attributs plutôt que
stockée redondamment.

\section{Architecture logicielle}

L'application suit une architecture en couches classique
(Figure~\ref{fig:architecture}) : \texttt{Controller} $\rightarrow$
\texttt{Service} $\rightarrow$ \texttt{Repository} $\rightarrow$
\texttt{Entity}, adossée à PostgreSQL pour la persistance et Redis pour les
besoins transverses (limitation de débit, verrous, état éphémère des
tentatives de quiz).

\begin{figure}[H]
  \centering
  \includegraphics[width=0.75\textwidth,height=0.85\textheight,keepaspectratio]{diagramme_architecture.png}
  \caption{Architecture en couches}
  \label{fig:architecture}
\end{figure}

La règle de conception appliquée de façon constante dans le code est la
suivante : un contrôleur ne fait que du routage HTTP, de la validation
\texttt{@PreAuthorize} et de la traduction requête/réponse ; toute la
logique métier --- y compris les vérifications de propriété de ressource
détaillées au chapitre~\ref{ch:realisation} --- vit dans la couche
service, jamais dans le contrôleur.

\section{Sécurité}

La sécurité de la plateforme ne pouvait pas être traitée comme une simple
case à cocher en fin de développement : dès lors qu'une partie du contenu
--- réponses de quiz, dossiers de stage, échanges de messagerie --- touche
directement à la vie privée et à l'évaluation d'un apprenant mineur ou
jeune majeur, l'authentification et le contrôle d'accès devaient être
pensés comme des contraintes de conception à part entière, et non comme
une couche ajoutée après coup. Cette section détaille les deux piliers sur
lesquels repose cette sécurité : l'authentification, qui répond à la
question \emph{qui êtes-vous ?}, et l'autorisation, qui répond à la
question, en apparence proche mais en réalité bien plus difficile à
garantir correctement, \emph{avez-vous le droit d'agir sur cette
ressource précise ?}

\subsection{Authentification}

L'authentification de la plateforme repose sur un jeton JWT (signé en
HS256) transporté par un cookie \texttt{httpOnly} --- c'est-à-dire un
cookie que le JavaScript exécuté dans le navigateur ne peut ni lire ni
manipuler, ce qui ferme la voie la plus commune d'exfiltration de jeton en
cas de faille de type \emph{cross-site scripting} (XSS) sur le frontend.
Le déroulement complet, depuis la soumission du formulaire de connexion
jusqu'à l'établissement du contexte de sécurité sur chaque requête
suivante, se décompose en quatre étapes. Premièrement,
\texttt{AuthService.login()} authentifie les identifiants soumis par
l'intermédiaire de l'\texttt{AuthenticationManager} de Spring Security,
qui compare le mot de passe fourni au hachage BCrypt stocké en base ---
jamais le mot de passe en clair. Deuxièmement, si l'utilisateur a activé
la double authentification par application TOTP (Google Authenticator ou
équivalent), la réponse renvoyée à ce stade n'est pas encore le cookie de
session final, mais un jeton temporaire valable cinq minutes, qui exige un
second appel explicite à l'endpoint \texttt{/verify-2fa} avant que la
connexion ne soit réellement effective --- une étape supplémentaire
délibérément conçue pour qu'un mot de passe compromis seul ne suffise
jamais à usurper un compte protégé par 2FA. Troisièmement, une fois
l'authentification (simple ou renforcée par 2FA) confirmée,
\texttt{JwtService} construit le jeton définitif, dont le sujet est
l'identifiant de l'utilisateur et dont les revendications embarquent son
adresse email et son rôle applicatif. Quatrièmement, à chaque requête
suivante, \texttt{JwtAuthenticationFilter} --- un filtre Spring exécuté
avant le filtre d'authentification standard du framework --- extrait ce
jeton du cookie, le valide, puis peuple le contexte de sécurité de la
requête courante avant que celle-ci n'atteigne le contrôleur ciblé.

Ce choix architectural a une conséquence directe sur la politique de
session retenue : le backend fonctionne en mode \texttt{STATELESS},
c'est-à-dire qu'aucune session n'est conservée en mémoire côté serveur.
L'intégralité de l'état d'authentification est reconstruite, à chaque
requête, à partir du seul contenu du jeton --- un choix cohérent avec une
architecture où le frontend Next.js est physiquement découplé du backend
Spring Boot, et qui simplifie d'autant la mise à l'échelle horizontale de
l'API si elle devait un jour être répliquée sur plusieurs instances,
puisqu'aucune session partagée n'a besoin d'être synchronisée entre elles.

\subsection{Autorisation}
\label{sec:autorisation}

Savoir \emph{qui} fait la requête ne suffit pas à savoir si cette requête
doit être acceptée : le contrôle d'accès de la plateforme combine, de
façon volontairement distincte, deux niveaux de vérification qui ne
répondent pas à la même question. Le premier niveau est un
\textbf{RBAC déclaratif}, appliqué au moyen de l'annotation
\texttt{@PreAuthorize} posée sur chaque endpoint sensible, et résolu
contre la hiérarchie de rôles \texttt{ROLE\_SUPER\_ADMIN > ROLE\_ADMIN}
présentée au chapitre~\ref{ch:realisation}. Ce niveau répond correctement
à la question \og cet utilisateur a-t-il le \emph{type} de compte requis
pour accéder à cette fonctionnalité ? \fg, mais il est structurellement
incapable de répondre à une seconde question, tout aussi critique dès lors
qu'une ressource appartient nommément à un utilisateur précis : \og cet
utilisateur a-t-il le droit d'accéder à \emph{cette instance précise} de
la ressource, et pas seulement à des ressources de ce type en général ?
\fg. C'est pour répondre à cette seconde question qu'un second niveau, un
\textbf{contrôle de propriété} implémenté au niveau de la couche service,
vient compléter le RBAC pour les ressources intrinsèquement personnelles
--- une conversation de messagerie, une notification, un dossier de stage
--- comme le détaille et l'illustre concrètement le
\S\ref{sec:acces-proprietaire} du chapitre suivant.

\section{Choix technologiques}

Chaque brique technologique retenue pour ce projet répond à une
contrainte identifiée dans les sections précédentes, plutôt qu'à un
simple effet de mode : le Tableau~\ref{tab:stack} met en regard chaque
choix et la justification qui le sous-tend, en particulier pour les
décisions les moins évidentes de prime abord --- pourquoi Redis en plus
d'une base relationnelle déjà présente, ou pourquoi un rendu serveur côté
frontend pour une seule page de la plateforme.

\begin{table}[H]
  \centering
  \caption{Stack technique et justification}
  \label{tab:stack}
  \begin{tabular}{@{}p{3.2cm}p{9cm}@{}}
    \toprule
    Technologie & Justification \\
    \midrule
    Spring Boot 3 / Java 21 & Écosystème mature, sécurité (Spring Security)
      et accès aux données (Spring Data JPA) intégrés, adapté à une
      application métier avec des règles de progression complexes. \\
    PostgreSQL + Flyway & Base relationnelle fiable pour un modèle fortement
      structuré (formations/modules/leçons) ; migrations versionnées
      reproductibles entre environnements. \\
    Redis & Structures de données en mémoire adaptées aux besoins
      transverses à courte durée de vie (compteurs de débit, verrous
      d'édition, état de tentative de quiz) sans alourdir PostgreSQL de
      lectures/écritures fréquentes et volatiles. \\
    JWT + cookie httpOnly & Authentification sans état, compatible avec un
      frontend découplé (Next.js), sans exposer le jeton au JavaScript
      (protection XSS). \\
    Next.js 15 (frontend) & Rendu serveur pour les pages publiques exigeant
      des métadonnées SEO/Open Graph (ex. vérification de certificat, voir
      \S\ref{sec:certificat-linkedin}). \\
    \bottomrule
  \end{tabular}
\end{table}

\section{Conclusion}

Ce chapitre a présenté les cas d'utilisation, le modèle de données, les
choix d'architecture et de sécurité qui structurent IAT Academy. Le
chapitre~\ref{ch:realisation} détaille maintenant comment ces choix se
traduisent dans l'implémentation, en particulier sur les mécanismes qui
dépassent le CRUD simple.
